\documentclass[10pt,a4paper]{amsart}
\usepackage{amsmath,amssymb,amsthm,graphicx}
\usepackage{verbatim}
\newenvironment{reference}{\comment}{\endcomment}
%\newenvironment{reference}{}{}
\newenvironment{slogan}{\comment}{\endcomment}
\newenvironment{history}{\comment}{\endcomment}

% For commutative diagrams we use Xy-pic
\usepackage[all]{xy}

% We use 2cell for 2-commutative diagrams.
\xyoption{2cell}
% Theorem environments.
%
\theoremstyle{plain}
\newtheorem{theorem}[subsection]{Theorem}
\newtheorem{proposition}[subsection]{Proposition}
\newtheorem{lemma}[subsection]{Lemma}

\theoremstyle{definition}
\newtheorem{definition}[subsection]{Definition}
\newtheorem{example}[subsection]{Example}
\newtheorem{exercise}[subsection]{Exercise}
\newtheorem{situation}[subsection]{Situation}

\theoremstyle{remark}
\newtheorem{remark}[subsection]{Remark}
\newtheorem{remarks}[subsection]{Remarks}

\numberwithin{equation}{subsection}

% Macros
%
\def\lim{\mathop{\mathrm{lim}}\nolimits}
\def\colim{\mathop{\mathrm{colim}}\nolimits}
\def\Spec{\mathop{\mathrm{Spec}}}
\def\Hom{\mathop{\mathrm{Hom}}\nolimits}
\def\Ext{\mathop{\mathrm{Ext}}\nolimits}
\def\SheafHom{\mathop{\mathcal{H}\!\mathit{om}}\nolimits}
\def\SheafExt{\mathop{\mathcal{E}\!\mathit{xt}}\nolimits}
\def\Sch{\mathit{Sch}}
\def\Mor{\mathop{\mathrm{Mor}}\nolimits}
\def\Ob{\mathop{\mathrm{Ob}}\nolimits}
\def\Sh{\mathop{\mathit{Sh}}\nolimits}
\def\NL{\mathop{N\!L}\nolimits}
\def\CH{\mathop{\mathrm{CH}}\nolimits}
\def\proetale{{pro\text{-}\acute{e}tale}}
\def\etale{{\acute{e}tale}}
\def\QCoh{\mathit{QCoh}}
\def\Ker{\mathop{\mathrm{Ker}}}
\def\Im{\mathop{\mathrm{Im}}}
\def\Coker{\mathop{\mathrm{Coker}}}
\def\Coim{\mathop{\mathrm{Coim}}}

% Boxtimes
%
\DeclareMathSymbol{\boxtimes}{\mathbin}{AMSa}{"02}

%
% Macros for moduli stacks/spaces
%
\def\QCohstack{\mathcal{QC}\!\mathit{oh}}
\def\Cohstack{\mathcal{C}\!\mathit{oh}}
\def\Spacesstack{\mathcal{S}\!\mathit{paces}}
\def\Quotfunctor{\mathrm{Quot}}
\def\Hilbfunctor{\mathrm{Hilb}}
\def\Curvesstack{\mathcal{C}\!\mathit{urves}}
\def\Polarizedstack{\mathcal{P}\!\mathit{olarized}}
\def\Complexesstack{\mathcal{C}\!\mathit{omplexes}}
% \Pic is the operator that assigns to X its picard group, usage \Pic(X)
% \Picardstack_{X/B} denotes the Picard stack of X over B
% \Picardfunctor_{X/B} denotes the Picard functor of X over B
\def\Pic{\mathop{\mathrm{Pic}}\nolimits}
\def\Picardstack{\mathcal{P}\!\mathit{ic}}
\def\Picardfunctor{\mathrm{Pic}}
\def\Deformationcategory{\mathcal{D}\!\mathit{ef}}
\usepackage{cleveref}
\crefname{originalreference}{original reference}{original references}
\Crefname{originalreference}{Original reference}{Original references}
\makeatletter
\newlabel{item-direct-sum}{{07Q1}{source}{}{sourcepointer-1}{}}
\newlabel{item-direct-sum@cref}{{[originalreference][0][]07Q1}{[1][0][]source}{}{}{}}
\makeatother
\crefname{originalreference}{original reference}{original references}
\Crefname{originalreference}{Original reference}{Original references}
\makeatletter
\newlabel{item-factor}{{07Q5}{source}{}{sourcepointer-2}{}}
\newlabel{item-factor@cref}{{[originalreference][0][]07Q5}{[1][0][]source}{}{}{}}
\makeatother
\crefname{originalreference}{original reference}{original references}
\Crefname{originalreference}{Original reference}{Original references}
\makeatletter
\newlabel{item-horizontal}{{07Q6}{source}{}{sourcepointer-3}{}}
\newlabel{item-horizontal@cref}{{[originalreference][0][]07Q6}{[1][0][]source}{}{}{}}
\makeatother
\crefname{originalreference}{original reference}{original references}
\Crefname{originalreference}{Original reference}{Original references}
\makeatletter
\newlabel{item-d-a-zero}{{07Q0}{source}{}{sourcepointer-4}{}}
\newlabel{item-d-a-zero@cref}{{[originalreference][0][]07Q0}{[1][0][]source}{}{}{}}
\makeatother
\crefname{originalreference}{original reference}{original references}
\Crefname{originalreference}{Original reference}{Original references}
\makeatletter
\newlabel{item-injective}{{07Q4}{source}{}{sourcepointer-5}{}}
\newlabel{item-injective@cref}{{[originalreference][0][]07Q4}{[1][0][]source}{}{}{}}
\makeatother
\crefname{originalreference}{original reference}{original references}
\Crefname{originalreference}{Original reference}{Original references}
\makeatletter
\newlabel{item-integrate}{{07Q3}{source}{}{sourcepointer-6}{}}
\newlabel{item-integrate@cref}{{[originalreference][0][]07Q3}{[1][0][]source}{}{}{}}
\makeatother
\title{Constructing the homotopy inverse to a base-change map of de Rham complexes}
\author{The Stacks Project Authors}
\date{stacks-project:crystalline.tex:local-snapshot-20261012; GNU FDL 1.2 or later}
\begin{document}
\maketitle
\noindent\textit{Editorial note.} All delivered mathematical text is the author's own, copied verbatim from the pinned Stacks Project chapter file crystalline.tex; no statement or proof step was rewritten, reordered or removed.\par\medskip
\noindent\textit{Editorial note.} Span recomputation by label pairing (LANE\_BRIEF 10.2): in this chapter the author's proof environment follows the statement environment as a standalone block; the spec covers the complete environment pair (statement span, proof span) as recorded in the pool candidate.\par\medskip
\ifdefined\hypertarget\hypertarget{sourcepointer-1}{}\fi
\noindent\textit{Source reference.} Original reference 07Q1: sources/stacks-chow/tags.txt, source lines 9816--9816. Cross-chapter standard prerequisite cited by the author inside the counted span; it is used as a pointer and its own proof is not part of this item.\par\medskip
\ifdefined\hypertarget\hypertarget{sourcepointer-2}{}\fi
\noindent\textit{Source reference.} Original reference 07Q5: sources/stacks-chow/tags.txt, source lines 9820--9820. Cross-chapter standard prerequisite cited by the author inside the counted span; it is used as a pointer and its own proof is not part of this item.\par\medskip
\ifdefined\hypertarget\hypertarget{sourcepointer-3}{}\fi
\noindent\textit{Source reference.} Original reference 07Q6: sources/stacks-chow/tags.txt, source lines 9821--9821. Cross-chapter standard prerequisite cited by the author inside the counted span; it is used as a pointer and its own proof is not part of this item.\par\medskip
\ifdefined\hypertarget\hypertarget{sourcepointer-4}{}\fi
\noindent\textit{Source reference.} Original reference 07Q0: sources/stacks-chow/tags.txt, source lines 9815--9815. Cross-chapter standard prerequisite cited by the author inside the counted span; it is used as a pointer and its own proof is not part of this item.\par\medskip
\ifdefined\hypertarget\hypertarget{sourcepointer-5}{}\fi
\noindent\textit{Source reference.} Original reference 07Q4: sources/stacks-chow/tags.txt, source lines 9819--9819. Cross-chapter standard prerequisite cited by the author inside the counted span; it is used as a pointer and its own proof is not part of this item.\par\medskip
\ifdefined\hypertarget\hypertarget{sourcepointer-6}{}\fi
\noindent\textit{Source reference.} Original reference 07Q3: sources/stacks-chow/tags.txt, source lines 9818--9818. Cross-chapter standard prerequisite cited by the author inside the counted span; it is used as a pointer and its own proof is not part of this item.\par\medskip
\begin{lemma}
\label{lemma-find-homotopy}
In the situation above there exists a map of complexes
$$
e_M^\bullet :
M \otimes_B (\Omega')^\bullet
\longrightarrow
M \otimes_B \Omega^\bullet
$$
such that $c_M^\bullet \circ e_M^\bullet$
and $e_M^\bullet \circ c_M^\bullet$ are homotopic to
multiplication by $a$.
\end{lemma}

\begin{proof}
In this proof all tensor products are over $B$.
Assumption (\ref{item-direct-sum}) implies that
$$
M \otimes (\Omega')^i =
(B' \otimes M \otimes \Omega^i)
\oplus
(B' \text{d}z \otimes M \otimes \Omega^{i - 1})
$$
for all $i \geq 0$. A collection of additive generators for
$M \otimes (\Omega')^i$ is formed by elements of the form
$f \omega$ and elements of the form $f \text{d}z \wedge \eta$
where $f \in B'$, $\omega \in M \otimes \Omega^i$, and
$\eta \in M \otimes \Omega^{i - 1}$.

\medskip\noindent
For $f \in B'$ we write
$$
\epsilon(f) = af - \theta(\partial_z(f))
\quad\text{and}\quad
\epsilon'(f) = (\theta \otimes 1)(\text{d}_1(f)) - \text{d}_1(\theta(f))
$$
so that $\epsilon(f) \in B$ and $\epsilon'(f) \in \Omega$ by
assumptions (\ref{item-factor}) and (\ref{item-horizontal}).
We define $e_M^\bullet$ by the rules
$e^i_M(f\omega) = \epsilon(f) \omega$ and
$e^i_M(f \text{d}z \wedge \eta) = \epsilon'(f) \wedge \eta$.
We will see below that the collection of maps $e^i_M$ is a map of complexes.

\medskip\noindent
We define
$$
h^i : M \otimes_B (\Omega')^i \longrightarrow M \otimes_B (\Omega')^{i - 1}
$$
by the rules $h^i(f \omega) = 0$ and
$h^i(f \text{d}z \wedge \eta) = \theta(f) \eta$
for elements as above.  We claim that
$$
\text{d} \circ h + h \circ \text{d} = a - c_M^\bullet \circ e_M^\bullet
$$
Note that multiplication by $a$ is a map of complexes
by (\ref{item-d-a-zero}). Hence, since $c_M^\bullet$ is an injective map
of complexes by assumption (\ref{item-injective}), we conclude that
$e_M^\bullet$ is a map of complexes. To prove the claim we compute
\begin{align*}
(\text{d} \circ h + h \circ \text{d})(f \omega)
& =
h\left(\text{d}(f) \wedge \omega + f \nabla(\omega)\right)
\\
& =
\theta(\partial_z(f)) \omega
\\
& =
a f\omega - \epsilon(f)\omega 
\\
& =
a f \omega - c^i_M(e^i_M(f\omega))
\end{align*}
The second equality because $\text{d}z$ does not occur in $\nabla(\omega)$
and the third equality by assumption (6). Similarly, we have
\begin{align*}
(\text{d} \circ h + h \circ \text{d})(f \text{d}z \wedge \eta)
& =
\text{d}(\theta(f) \eta) +
h\left(\text{d}(f) \wedge \text{d}z \wedge \eta -
f \text{d}z \wedge \nabla(\eta)\right)
\\
& =
\text{d}(\theta(f)) \wedge \eta + \theta(f) \nabla(\eta)
- (\theta \otimes 1)(\text{d}_1(f)) \wedge \eta
- \theta(f) \nabla(\eta)
\\
& =
\text{d}_1(\theta(f)) \wedge \eta +
\partial_z(\theta(f)) \text{d}z \wedge \eta -
(\theta \otimes 1)(\text{d}_1(f)) \wedge \eta
\\
& =
a f \text{d}z \wedge \eta - \epsilon'(f) \wedge \eta \\
& = a f \text{d}z \wedge \eta - c^i_M(e^i_M(f \text{d}z \wedge \eta))
\end{align*}
The second equality because
$\text{d}(f) \wedge \text{d}z \wedge \eta =
- \text{d}z \wedge \text{d}_1(f) \wedge \eta$.
The fourth equality by assumption (\ref{item-integrate}).
On the other hand it is immediate from the definitions
that $e^i_M(c^i_M(\omega)) = \epsilon(1) \omega = a \omega$.
This proves the lemma.
\end{proof}

\end{document}
